\section{Conclusion and future scope}\label{sec:conclusion}

\paragraph{What was shown.} When tabular data contain repeated rows, the usual random-split check of conformal prediction can report 90\% coverage while new cases are covered much less often: \NCsMcOneNew{}\% on NASA mc1. A check that needs no model flagged every fresh data set with a real shortfall (\NScpHit{} of \NScpPos{}). Calibrating seen and new cases separately kept new-case coverage near the target on all \NScpN{} fresh data sets, at almost no cost. Most standard benchmarks are not affected, and the check says so.

\paragraph{What was not shown.} The predicted gap is a limit and a screen, not a forecast of the exact shortfall. The fix promises coverage over all new cases, not over each class. Exact matching misses near-copies. All evidence comes from random splits of public data sets, not from a deployed system.

\paragraph{Next steps.}
\begin{enumerate}
\item Extend the check from exact to near-identical records with a principled choice of matching resolution.
\item Combine the seen/new split with class-wise calibration for rare classes.
\item Test the routine on time-ordered deployments, where future records are new by construction.
\item Integrate the check into common conformal libraries, so that new-case coverage is reported by default.
\end{enumerate}
